\section{Implementación}\label{cap:implementacion}

Esta sección describe cómo se materializan en código las decisiones de diseño de la sección~\ref{cap:arquitectura}. La narrativa sigue el orden lógico de la plataforma: fuente técnica, ingesta, núcleo Python, configuración funcional, consola web y, por último, el despliegue reproducible. La organización del repositorio se describe en la sección~\ref{sec:metodologia-repo}. Los comandos exactos para reproducir cada bloque están en el Anexo~\ref{ap:reproduccion}; las capturas que evidencian cada paso y los listados de configuración extensos, en el Anexo~\ref{ap:capturas}; y el recorrido visual completo desde la consola web, en el Anexo~\ref{ap:demo-operativa}.\ct{El tutor indica que esto es Resultados, no ``Implementación'': en la estructura normativa la implementación forma parte de Resultados.}

\subsection{Fuente técnica reproducible: PostgreSQL en tres capas}\label{sec:implementacion-postgres}

La fuente técnica es un PostgreSQL desplegado en el mismo clúster Kubernetes que \ac{OM}, inicializado a partir de \texttt{sql/opendata\_demo\_init.sql}~\cite{postgresqlDocs}.\ct{Hecho: la arquitectura precede ahora a la implementación dentro de Resultados, como pedía el tutor.} El \ac{SQL} crea tres esquemas alineados con el patrón \emph{medallion} (bronze/silver/gold), de modo que el caso de uso de validación dispone de datos sintéticos verosímiles sin depender de descargas externas. Únicamente el esquema \texttt{gold} se considera publicable: las capas \texttt{bronze} y \texttt{silver} se ingieren para mantener visible la cadena de transformaciones, pero no se exportan como datasets \ac{DCATAPES}, por las razones de conformidad semántica, alineamiento con \ac{DMBOK} y separación entre catálogo técnico y de publicación que se argumentan en la sección~\ref{sec:decision-gold-only}. El listado~\ref{code:postgres-gold} muestra la creación de una de las tablas \texttt{gold} publicables, con una estructura suficientemente expresiva (claves temporales, dimensiones territoriales y métricas) para que el catálogo exportado contenga campos significativos al validar con \ac{SHACL}.

\begin{lstlisting}[language=SQL, caption={Fragmento de \texttt{sql/opendata\_demo\_init.sql}: tabla publicable de la capa gold.}, label=code:postgres-gold, captionpos=b]
CREATE TABLE IF NOT EXISTS gold.agenda_cultural_publica (
    evento_id        BIGINT      PRIMARY KEY,
    fecha            DATE        NOT NULL,
    municipio        TEXT        NOT NULL,
    categoria        TEXT        NOT NULL,
    titulo           TEXT        NOT NULL,
    asistencia_est   INTEGER     CHECK (asistencia_est >= 0),
    actualizado_en   TIMESTAMPTZ DEFAULT now()
);
\end{lstlisting}

\subsection{Doble ingesta técnica}\label{sec:implementacion-ingesta}

En OpenMetadata, una \emph{ingesta} consiste en conectarse a una fuente de datos y registrar automáticamente sus metadatos técnicos (el servicio, las bases de datos, los esquemas, las tablas y las columnas), sin copiar los datos de negocio. La \emph{doble ingesta} ejecuta ese proceso dos veces sobre la misma base PostgreSQL, registrándola bajo dos servicios lógicos distintos en \ac{OM}: \fqn{postgres_demo_service} y \fqn{postgres_validation_service}.\ct{El tutor avisa de que la doble ingesta no se entiende para quien la lea sin contexto. Se explica qué es una ingesta en OpenMetadata, por qué se ejecuta dos veces y qué problema (tablas homónimas en fuentes distintas) sirve para validar.}

El motivo no es duplicar datos, sino reproducir de forma controlada un escenario habitual en organizaciones reales: que una misma tabla, por ejemplo \fqn{gold.agenda_cultural_publica}, exista con idéntico nombre en dos fuentes diferentes. Con ello se comprueba que la lógica de gobierno identifica cada activo por su nombre completo en \ac{OM} (\ac{FQN}, que incluye servicio, base de datos, esquema y tabla) y no por el nombre corto de la tabla. De lo contrario, las dos tablas homónimas se mezclarían y el catálogo resultante sería incorrecto. La ingesta se invoca desde \fqn{scripts/infra/ingest_postgres_double.ps1} y su resultado es visible en la \ac{UI} de \ac{OM} (figura~\ref{fig:om-services}): cuatro tablas publicables, dos por servicio, cada una con su \ac{FQN} único.

\subsection{Núcleo Python: paquete \texttt{om\_dcat\_sync}}\label{sec:implementacion-nucleo}

El paquete \texttt{tfm\_ingestor} (alias importable \texttt{om\_dcat\_sync}) concentra toda la lógica de gobierno, exportación y validación. Está estructurado en módulos con responsabilidades acotadas, conforme al principio de núcleo único declarado en la sección~\ref{cap:arquitectura}. La tabla~\ref{tab:modulos-python} lista los módulos principales y la figura~\ref{fig:modulos-python} del Anexo~\ref{ap:capturas} resume su composición.\ct{El tutor considera la figura de la estructura interna del paquete candidata a moverse a un anexo.}

\begin{table}[h]
\footnotesize
\centering
\begin{tabular}{|L{4cm}|L{8cm}|}
\hline
\rowcolor{tema!10}
\bft{Módulo} & \bft{Responsabilidad} \\ \hline
\texttt{workflow\_service.py} & Orquestador canónico del workflow end-to-end. \\ \hline
\texttt{governance\_service.py} & Sincroniza gobierno entre hoja \ac{CSV} y \ac{OM}. \\ \hline
\texttt{governance\_sheet.py} & Lectura, escritura y validación de la hoja \ac{CSV}. \\ \hline
\texttt{om\_api.py} & Cliente \ac{REST} ligero para la \ac{API} de \ac{OM}. \\ \hline
\texttt{dcat\_export.py} & Construye el grafo \ac{DCATAPES} y lo serializa en \ac{JSONLD}. \\ \hline
\texttt{shacl\_validation.py} & Aplica las \emph{shapes} vendorizadas con pySHACL. \\ \hline
\texttt{config.py} & Carga defaults \ac{YAML} y reglas de mapeo. \\ \hline
\texttt{operational\_profile.py} & Une defaults, reglas y caso de validación. \\ \hline
\texttt{runtime\_validation.py} & Verifica coherencia del estado vivo de \ac{OM}. \\ \hline
\texttt{main.py} & \ac{CLI} basado en \texttt{argparse}. \\ \hline
\end{tabular}
\caption{Módulos principales de \texttt{tfm\_ingestor}. Elaboración propia.}
\label{tab:modulos-python}
\end{table}

\subsubsection{Configuración inmutable y entrypoint}

El módulo \fqn{workflow_service.py} define una \texttt{dataclass} inmutable (\texttt{frozen=True})~\cite{pythonDataclasses} que fija todos los parámetros del workflow antes de su ejecución, lo que evita estados mutables compartidos entre llamadas. El entrypoint \texttt{run\_workflow} encadena las cuatro etapas (refresco de hoja, sincronización de gobierno, exportación y validación); su esqueleto completo, que amplía el listado~\ref{code:workflow-service} con la secuencia operativa real, se recoge en el listado~\ref{code:workflow-service-impl} del Anexo~\ref{ap:capturas}.

\subsubsection{Sincronización repetible (idempotente) de gobierno}

En este contexto, \emph{sincronización idempotente} significa, sencillamente, que la operación es \textbf{repetible sin efectos colaterales}: ejecutarla una vez o varias veces seguidas deja OpenMetadata en el mismo estado. La primera ejecución aplica los cambios necesarios y las siguientes no vuelven a modificar nada.\ct{El tutor sigue sin ver clara la palabra ``idempotente'' en este contexto. Se renombra el subtítulo a ``Sincronización repetible (idempotente)'' y se abre con una definición llana: repetible sin efectos colaterales.} \texttt{governance\_service.py} mapea cada fila de la hoja \ac{CSV} a operaciones \texttt{PATCH} o \texttt{POST} sobre la entidad correspondiente en \ac{OM} usando su \ac{FQN}. Esa propiedad se materializa por dos vías: las propiedades gestionadas (\texttt{MANAGED\_DCAT\_CUSTOM\_PROPERTIES}) solo se reescriben si cambian, y las etiquetas \texttt{dcat\_theme.*} se aplican como conjuntos, nunca como acumulación lineal. La consecuencia es que ejecutar el workflow dos veces consecutivas produce cero cambios en la segunda ejecución, una de las métricas de validación de la sección~\ref{cap:validacion}.

El resultado del gobierno es directamente observable en \ac{OM}: tras el workflow, cada tabla \texttt{gold} publicable conserva sus propiedades personalizadas DCAT (\fqn{dcat_publisher_name}, \fqn{dcat_hvd_category}, \fqn{dcat_access_url}) y sus etiquetas \fqn{dcat_theme.*}, que son la materia prima que el exportador transforma en el grafo \ac{DCATAPES} (figura~\ref{fig:om-tabla-propiedades}, Anexo~\ref{ap:capturas}).

\subsubsection{Exportador \texorpdfstring{\ac{DCATAPES}}{DCAT-AP-ES}}

\texttt{dcat\_export.py} compone el grafo conceptual a partir de las clases activas \ident{dcat:Catalog}, \ident{dcat:Dataset}, \ident{dcat:Distribution}, \ident{dcat:DataService} y \ident{foaf:Agent}. La construcción se hace en Python puro mediante diccionarios anidados que respetan el contexto \ac{JSONLD}, y cuando es necesario manipular el grafo (normalización o deduplicación) se delega en RDFLib~\cite{rdflib}. El contexto \ac{JSONLD} canónico, con prefijos estables, idioma por defecto en español y enlace a los vocabularios temáticos de NTI-RISP~\cite{datosGobNtiRisp}, se recoge en el listado~\ref{code:dcat-export-context} del Anexo~\ref{ap:capturas}.

\subsubsection{Validación SHACL del catálogo exportado}

Esta etapa comprueba formalmente que el catálogo cumple el perfil DCAT-AP-ES. \fqn{shacl_validation.py} toma como entrada el grafo \ac{JSONLD} exportado en la etapa anterior y un \emph{grafo de shapes} (el conjunto de restricciones del perfil), y los confronta con el motor pySHACL~\cite{pyshacl}, implementación de referencia de \ac{SHACL} para Python. El concepto de \emph{shape} y un ejemplo mínimo se introdujeron en la sección~\ref{sec:shacl}. Aquí se aplican las shapes oficiales completas.\ct{El tutor pide desarrollar más esta sección. Se amplía con el flujo de validación (entrada JSON-LD y grafo de shapes, motor pySHACL), el papel del caso HVD, la estructura del informe, la distinción Violation/Warning y el criterio de aceptación.}

El bundle de shapes está vendorizado en \fqn{tfm_ingestor/resources/shacl/} y congelado en el commit \texttt{f2c8a888} del repositorio oficial \fqn{datosgobes/DCAT-AP-ES}~\cite{datosGobValidacionDcatApEs2025}. Congelarlo, en lugar de descargarlo en cada ejecución, evita dependencias de red y garantiza que dos validaciones del mismo catálogo den siempre el mismo resultado. El caso \texttt{hvd} carga, además de las restricciones generales, las shapes específicas de los conjuntos de datos de alto valor, que exigen propiedades adicionales como la categoría HVD, la legislación aplicable y el servicio de acceso.

El resultado de pySHACL tiene dos partes: un valor booleano \texttt{conforms}, que indica si el grafo cumple todas las restricciones \emph{obligatorias}, y un informe en Turtle que detalla cada incidencia mediante \texttt{sh:focusNode} (el recurso afectado), \texttt{sh:resultPath} (la propiedad implicada), \texttt{sh:severity} (la gravedad) y \texttt{sh:resultMessage} (la explicación). La gravedad distingue dos casos con tratamiento distinto. Una \texttt{sh:Violation} señala el incumplimiento de una restricción obligatoria y hace que \texttt{conforms} sea falso. Un \texttt{sh:Warning} indica solo una recomendación no cubierta y no invalida la conformidad. Por eso el módulo postprocesa el informe para filtrar los \texttt{sh:Warning} cuando se ejecuta con \texttt{-{}-allow-warnings}, dejando visibles únicamente las incidencias bloqueantes. El criterio de aceptación es, por tanto, inequívoco: el catálogo se considera conforme cuando no contiene ninguna \texttt{sh:Violation}. Este resultado se consolida como la métrica de conformidad SHACL de la sección~\ref{cap:validacion} y se contrasta con un validador europeo independiente en la sección~\ref{sec:validacion-europea}.

\subsection{Hoja funcional de gobierno y defaults globales}\label{sec:implementacion-hoja}

La curación funcional se concentra en \texttt{tfm\_ingestor/config/gold\_governance.csv}, cuya estructura columnar se describe en la tabla~\ref{tab:hoja-gobierno-columnas} y se detalla con criterios de validación en el Anexo~\ref{ap:ingesta-gobierno}. El uso de \ac{CSV} con separador \texttt{;} no es accidental: facilita la edición en hojas de cálculo locales sin Excel y, sobre todo, hace que el diff en Git sea legible para revisión editorial. Los defaults globales viven en \texttt{governance\_defaults.yaml}, que centraliza publicador, homepage, idioma, licencia, contacto y la legislación \ac{HVD} aplicable. La pantalla \emph{Gobierno} de la consola web (figura~\ref{fig:web-pantalla-gobierno}, Anexo~\ref{ap:capturas}) actúa sobre la misma hoja y aplica las mismas reglas de validación que el \ac{CLI}; su uso real se narra en la sección~\ref{sec:demo-gobierno} del Anexo~\ref{ap:demo-operativa}.

\subsection{Consola web operativa Next.js}\label{sec:implementacion-web}

La consola web se construye con Next.js~\cite{nextjsDocs}, React~\cite{reactDocs} y TypeScript~\cite{typescriptHandbook}, organizada en rutas que mapean uno a uno con las etapas del workflow: \emph{Infraestructura}, \emph{Ingesta}, \emph{Gobierno}, \emph{Workflow}, \emph{DCAT}, \emph{Estado vivo}, \emph{Validación}, \emph{SHACL}, \emph{Artefactos} y \emph{Ejecuciones}. La arquitectura interna respeta la regla de núcleo único: la \ac{UI} nunca habla directamente con \ac{OM} ni reimplementa reglas, sino que invoca una lista cerrada de operaciones de servidor (\texttt{operations.ts}) que ejecutan el \ac{CLI} o los scripts PowerShell. La descripción pantalla a pantalla, con los artefactos producidos en cada paso y sus capturas, se desarrolla en el Anexo~\ref{ap:demo-operativa}, y los comandos equivalentes desde \ac{CLI} y \texttt{pnpm} en el Anexo~\ref{ap:validacion-web}.

Cada ejecución crea un \emph{job} versionado en \texttt{state/web\_jobs/} con identificador único, log completo, código de salida, duración y artefactos generados. El listado~\ref{code:web-job-state} del Anexo~\ref{ap:capturas} muestra la estructura de un \emph{job} persistido, que la pantalla \emph{Ejecuciones} consume para renderizar el historial.

\subsection{Despliegue reproducible en Kubernetes}\label{sec:implementacion-k8s}

El despliegue local utiliza Kind como motor de clúster Kubernetes encima de Docker~\cite{dockerOverview, kindDocs, kubernetesConcepts}. \ac{OM} y sus dependencias (MySQL y OpenSearch) se instalan mediante Helm~\cite{helmDocs} usando los \emph{values} canónicos del repositorio, cuyo fragmento mínimo se recoge en el listado~\ref{code:k8s-values} del Anexo~\ref{ap:capturas}.\ct{Hecho: el tutor pedía explicar primero los módulos y dejar el despliegue de infraestructura para el final. Se ha movido esta subsección al cierre de la implementación.} El script \texttt{scripts/infra/launch\_infra.ps1} encadena la creación del clúster, la instalación de los charts y los \texttt{port-forward} necesarios para exponer la \ac{UI} de \ac{OM} en \texttt{localhost:8585} y el PostgreSQL de referencia en \texttt{localhost:5432}. La idempotencia se garantiza mediante \texttt{helm upgrade --install} y comprobaciones con \texttt{kubectl get} previas a cada paso destructivo. Los contenedores Docker y los \emph{pods} resultantes (figuras~\ref{fig:docker-containers} y~\ref{fig:pods-kubernetes}, Anexo~\ref{ap:capturas}), así como los \emph{values} Helm canónicos, se detallan en el Anexo~\ref{ap:infraestructura}.

\subsection{Resumen de implementación}

La implementación cumple las tres reglas duras declaradas en la sección~\ref{cap:arquitectura}: una única lógica de negocio en Python, idempotencia comprobable y separación estricta entre datos técnicos, curación funcional y derivaciones automáticas. La sección siguiente describe cómo se valida formalmente este ensamblaje y qué evidencias produce.
